\input format.tex

%\usepackage{arydshln}

\begin{document}

%%%
\begin{spacing}{1.25}

\setlength{\arrayrulewidth}{0.5pt}
\color{gray2}
\vspace*{8mm}
\fontsize{9.3pt}{9.8pt}\selectfont
\begin{tctabularx}{tabularx={p{15.4cm}<{\centering}p{0pt}}}
\\[-6.769pt]
\cellcolor{topcolor} \raisebox{1pt}{\color{gray0}\xiaosihao\lantxh\bf 益生菌调理方案} \\
\end{tctabularx}

{
\lantxh
\vspace*{-4.25mm}
\noindent
\begin{longtable}{|p{2.25cm}|p{12.732cm}|}
\hline
\hspace*{5mm}\parbox[c]{\hsize}{\vskip10pt\raisebox{-1.5ex}{\includegraphics[scale=1.25]{probiotics_fangan.pdf}}\\{\xiaosihao\bf 益生菌}\vskip10pt} &
\parbox[c]{\hsize}{\vskip10pt\wuhao {建议每天适量补充含双歧杆菌、乳酸杆菌的复合益生菌固体饮料，可以调节肠道PH值，抑制肠道有害菌的生长，调理肠道微生态的平衡，能帮助减少肠道对糖类和脂肪的消化吸收，有利于控制体重。饭后1小时食用效果更佳。}\vskip10pt}\\
\hline
\end{longtable}
}

\vspace*{8mm}
\fontsize{9.3pt}{9.8pt}\selectfont
\begin{tctabularx}{tabularx={p{15.4cm}<{\centering}p{0pt}}}
\\[-6.769pt]
\cellcolor{topcolor} \raisebox{1pt}{\color{gray0}\xiaosihao\lantxh\bf 益生元调理方案} \\
\end{tctabularx}

{
\lantxh
\vspace*{-4.25mm}
\noindent
\begin{longtable}{|p{2.25cm}|p{12.732cm}|}
\hline
\hspace*{5mm}\parbox[c]{\hsize}{\vskip10pt\raisebox{-1.5ex}{\includegraphics[scale=1.25]{prebiotics_fangan.pdf}}\\{\xiaosihao\bf 益生元}\vskip10pt} &
\parbox[c]{\hsize}{\vskip10pt\wuhao {1.建议每天补充适量的可溶性膳食纤维，如聚葡萄糖。膳食纤维不能被人体直接消化吸收，可减少人体能量摄入，减少肠道对糖类和脂肪的消化吸收。
2.建议补充低聚糖类益生元产品，可以促进排泄物的排出，减少有害物质在体内的存留时间。如患炎症性肠病、慢性腹泻，不宜服用菊粉，建议选择低聚半乳糖及低聚果糖。
3.可补充含茶多酚、黄精提取物的功能食品，能有效帮助控制体重。}\vskip10pt}\\
\hline
\end{longtable}
}

\end{spacing}
\end{document}
